\subsubsection{Limits of the base potentials}
\label{sec:warped-endpoint-potential-limits}

The normalized potentials used to construct the interior balanced metrics have a unique continuous limit at each finite face. The compactness and uniqueness needed for this conclusion concern the closed base equation.

\Needspace{15\baselineskip}
\begin{proposition}\label{prop:warped-endpoint-potential-limits}\index{Gauduchon metric!boundary potential limits}\index{Compactness!warped endpoint potentials}
Under the hypotheses of \cref{thm:warped-unit-volume-completion}, let $\theta_j=b_j\kappa+L_jw\in\mathscr S$ converge in coefficients to a finite face $\theta_0=b\kappa+L_0w$. Normalize the smooth base potentials in \cref{thm:warped-gauduchon-volume} by $\sup_B\phi_j=0$ and put $u_j=\phi_j\circ p_B$. Then $u_j$ converge uniformly on $X$ to a continuous pulled-back solution $u_0$ of \eqref{eq:warped-continuous-endpoint-equation}. For every $1\leq k\leq n$,
\[
 (\theta_j+\ddc u_j)^k\rightharpoonup
             (\theta_0+\ddc u_0)^k.
\]
It is the unique normalized bounded solution pulled back from the base.
\end{proposition}
\begin{proof}
Write $v_j=c_j/e_j$ and $\psi_j=\phi_j/b_j$. The normalized closed base equation is
\[
 (\kappa+\ddc\psi_j)^m=F_{v_j}\sigma,
 \qquad F_v=\frac1{Q(v)(1+va)},\qquad\sup_B\psi_j=0.
\]
At the first face, $v_j\to\infty$. Set $\epsilon=1/v$ and $H_\epsilon=\int_B\sigma/(a+\epsilon)$. Then $F_v=[H_\epsilon(a+\epsilon)]^{-1}$ and $H_\epsilon\to H_a$. At the second face, $v_j\to-1/M$. With $f=1-a/M$ and $t=v+1/M>0$, the denominator is $Q(v)(f+ta)$ and $Q(v)\to H_-$. Thus the respective limits are
\[
 F_+=\frac1{H_a a},\qquad F_-=\frac1{H_-(1-a/M)}.
\]
Choose $1<p<r$, with $r>1$ the ratio for that face. Tubular integration as in \cref{prop:warped-continuous-endpoint-equation} gives $a^{-1}\in L^p(B,\sigma)$ or $f^{-1}\in L^p(B,\sigma)$, respectively. Near the first end, $F_v\leq C/a$; near the second, $F_v\leq C/f$. The extreme sets have smooth volume zero. Dominated convergence therefore gives $F_{v_j}\to F_\pm$ in $L^p(B,\sigma)$, and each limiting density has integral one. Since $\sigma/\kappa^m$ is smooth and strictly positive on compact $B$, the equation densities also have a common $L^p(B,\kappa^m)$ bound.

Apply \cref{thm:lp-uniform-compactness} to the fixed background $\kappa$. Every subsequence of $\psi_j$ has a uniformly convergent further subsequence. Bounded-product continuity identifies its limit with a normalized solution of $(\kappa+\ddc\psi_0)^m=F_\pm\sigma$. Closedness of $\kappa$ gives the Guan--Li condition, so \cref{uni-thm:GL-intro} makes this solution unique. Every cluster limit is consequently the same, which proves whole-sequence uniform convergence. Since $b_j\to b>0$, we obtain $u_j\to u_0=(b\psi_0)\circ p_B$ uniformly.

On each coordinate ball, coefficient convergence of the backgrounds gives a common smooth quadratic correction making all $u_j+q$ plurisubharmonic. These corrected potentials are uniformly bounded and converge uniformly. The capacity-continuity statement in \cref{sec:psh-local-capacity} gives weak convergence of all their bounded products, with uniform local mass bounds. Expanding off the correction as in \eqref{eq:bounded-expansion}, and using smooth convergence of the remaining background factors, gives the displayed current convergence. In top degree the interior products all equal $\mu$, so the limit solves the endpoint equation.

For any bounded pulled-back solution $v=\chi\circ p_B$ with $\sup_B\chi=0$, the horizontal current $A_\chi=b\kappa+\ddc\chi$ on $B$ is positive: pair its pulled-back current with positive base tests and $w$. The product expansion gives $nL_0A_\chi^m=\sigma$. Off the extreme set this determines $A_\chi^m=\sigma/(nL_0)$. That density already has total mass $H_0/n=b^m$, while closed bounded mass invariance gives $\int_BA_\chi^m=b^m$. Positivity excludes any residual mass on the extreme set. Thus $\chi/b$ satisfies the same normalized closed-base equation, and \cref{uni-thm:GL-intro} proves the asserted uniqueness.
\end{proof}

The statement selects the limit of the constructed base potentials. It does not identify arbitrary endpoint solutions or give a pointwise limit of the balanced metrics $\gamma_j=L_j^{-1/m}(b_j\kappa+\ddc\phi_j+L_jw)$.
